\documentclass{article}
\usepackage{iclr2026_conference,times}

\input{math_commands.tex}

\usepackage{booktabs}
\usepackage{hyperref}
\usepackage{url}
\usepackage{xcolor}

\title{Primitive-Constrained Verification and Repair for LLM-Generated Rigid-Body Animations}

\author{Anonymous Authors\\
Anonymous Institution\\
\texttt{anonymous@example.com}
}

%\iclrfinalcopy
\begin{document}

\maketitle

\begin{abstract}
Large language models can synthesize executable Blender programs from natural language, but direct keyframe generation often produces rigid-body animations that are visually plausible in isolated frames while violating basic contact constraints over time. Objects float above supports, penetrate tables or conveyors, detach from grippers, or satisfy a final language relation only by crossing through intervening geometry. We present a verification-oriented harness for LLM-generated rigid-body animation. The harness decomposes generation into a static scene stage and an animation extension stage, represents motion intent in an intermediate representation with stable object identities and contact constraints, executes Blender programs, verifies geometry with deterministic contact and penetration checks, and repairs common failures through measured scene-graph feedback. The current implementation further adds primitive-constrained animation helpers for support sliding, support sequences, pick-carry-place motion, pushing, dropping, and hinged rotation, together with auditable material texture resolution through online search, vision approval, and fallback. The animation helpers compute keyframes from evaluated world-space bounding boxes rather than prompt-level object centers, giving both generation and repair access to explicit contact geometry. Preliminary smoke tests verify the contact calculations used by these primitives on representative tabletop motions. We outline a controlled benchmark and ablation protocol for measuring validator-measured contact validity and refinement efficiency against unconstrained LLM keyframing.
\end{abstract}

\section{Introduction}

Language-conditioned 3D animation is becoming a practical interface for procedural content creation, robotics visualization, and simulation prototyping. A user should be able to request a box sliding across a table, a gripper transferring a cube into a tray, or a vehicle crossing a bridge, and receive an executable animation rather than a static asset or a textual plan. Recent large language models already write nontrivial graphics and tool-use programs \citep{chen2021evaluating,austin2021program,yao2023react,schick2023toolformer}, making Blender Python a natural target for this setting. The target requirement, however, is stronger than syntactic execution: an animation should preserve rigid-body contact, avoid penetration, keep relevant objects visible, and end in the intended relation.

In our development runs, direct LLM keyframing often fails because rigid-body animation correctness is governed by continuous geometry that is only weakly represented in text. A generated program may put a cube at the center of a table rather than at the table top plus half the cube height, may keep a carried object on the intended path while the gripper remains detached, or may satisfy start and end frames while interpolating through an obstacle. These errors suggest a mismatch between language-model token prediction and geometric invariants such as support footprint overlap, vertical contact gap, pairwise nonpenetration, and contact consistency across a frame window.

We address this mismatch by making LLM-generated animation verifier-facing before it is renderer-facing. Our system first asks a planner to produce a scene-and-motion intermediate representation (IR) with stable object ids, collision proxies, sampled frames, visibility requirements, and explicit contact constraints. A coder then emits a Blender script in two stages: a static scene is generated and verified first, and animation is added only after the scene baseline passes. Blender execution returns scene graphs, sampled transforms, screenshots, and animation previews. Deterministic validators check support, penetration, final placement, and carry-contact constraints over event boundaries and dense frame windows, while optional visual and video verifiers judge failures that geometry alone may miss.

The key technical idea is to combine verifier-facing constraints with primitive-constrained generation. Instead of asking the LLM to invent every keyframe, we expose reusable rigid-motion primitives whose semantics are expressed in measured world-space geometry: support sliding aligns an aggregate object bottom to a support top; support sequences place objects on successive surfaces at specified windows; pick-carry-place derives lift, carry, and release heights from source and destination supports; gripper spacing moves finger meshes outside the carried object's bounding box. When a script still violates contact constraints, targeted repair rewrites only the relevant motion path from measured scene-graph evidence.

This paper contributes an implementation-centered formulation of verifier-facing LLM animation generation:
\begin{itemize}
    \item a progressive scene-to-animation IR that makes object identity, contact windows, material texture policy, and verifier evidence explicit;
    \item a Blender execution harness that combines deterministic geometry checks, multi-view visual/video verification, and bounded refinement;
    \item primitive-constrained animation helpers that compute support, carry, push, drop, and hinge motions from measured world-space geometry; and
    \item local repair routines that rewrite support and carry failures using scene-graph measurements rather than unconstrained script regeneration.
\end{itemize}
We also document supporting engineering components, including vision-approved texture search, safe material fallback, media diagnostics, run traceability, and cache-friendly repeated coding prompts. Figure~\ref{fig:teaser} and Figure~\ref{fig:pipeline} show the intended qualitative evidence and pipeline structure.

\begin{figure}[t]
\centering
\fbox{\begin{minipage}[c][1.35in][c]{0.92\linewidth}
\centering
\textbf{Placeholder: teaser comparison.}\\
Left: direct LLM keyframes with floating, penetration, or detached carry contact.\\
Right: primitive-constrained animation with measured support/contact alignment.
\end{minipage}}
\caption{Qualitative teaser placeholder. The final figure should show representative failure modes and the corrected primitive-constrained outputs.}
\label{fig:teaser}
\end{figure}

\section{Related Work}

\paragraph{Language models as program generators.}
Large language models have been studied as code generators and tool-using agents \citep{chen2021evaluating,austin2021program,yao2023react,schick2023toolformer}. Blender scripting is an appealing target because scenes, materials, cameras, and keyframes can all be represented as executable programs. Our setting differs from general program synthesis because success depends not only on running code, but also on spatial and temporal constraints that must be measured in the executed 3D environment.

\paragraph{Text-to-3D and generative scene synthesis.}
Text-conditioned 3D generation methods synthesize shapes, assets, or scenes from natural language \citep{poole2022dreamfusion,lin2023magic3d}. These methods typically emphasize object appearance, asset quality, and scene-level visual fidelity. We instead study executable rigid-body animation programs with explicit object identities, support relations, contact windows, and verifier-readable failure reports.

\paragraph{Language-conditioned planning and motion primitives.}
Embodied and robotics systems often combine language models with structured plans, affordances, or motion primitives \citep{huang2022inner,ahn2022can}. We take inspiration from this decomposition, but our output is an offline graphics animation rather than a robot policy. The primitives in our system are kinematic and verification-oriented: they do not solve full dynamics, but they encode contact geometry that can be checked and repaired inside Blender.

\paragraph{Physical plausibility in animation.}
Physics engines can simulate rigid bodies, collisions, and constraints, but LLM-generated Blender scripts frequently bypass simulation and insert keyframes directly. Our harness adds a verifier-and-repair layer around keyframed animation. This is complementary to physics simulation: deterministic geometry checks can reject visible failures even when no simulator is used, and future versions can replace selected bounding-box checks with richer collision proxies or simulation-based validation.

\section{Method}

\subsection{Overview}

Our harness follows a generate--verify--repair loop around executable Blender programs. Given a language prompt, the planner emits a structured IR. The coder converts the IR into Blender Python. The runtime executes the script through a Blender addon server, records scene graphs and transforms, renders inspection views, and optionally renders GIF/MP4 previews. The verifier stack then combines static code checks, deterministic geometry checks, visual model feedback, and video model feedback. If any enabled verifier fails, the refiner receives structured reports and rendered evidence, updates the script, and the loop repeats under bounded refinement budgets.

\begin{figure}[t]
\centering
\fbox{\begin{minipage}[c][1.65in][c]{0.92\linewidth}
\centering
\textbf{Placeholder: pipeline figure.}\\
Prompt $\rightarrow$ scene/animation IR $\rightarrow$ static scene generation $\rightarrow$ deterministic + visual verification\\
$\rightarrow$ animation extension $\rightarrow$ contact/video verification $\rightarrow$ primitive repair/refinement.
\end{minipage}}
\caption{Pipeline placeholder. The final figure should emphasize that animation is added after a verified static scene baseline, and that deterministic evidence is shared by verification and repair.}
\label{fig:pipeline}
\end{figure}

\subsection{Verifier-Facing Intermediate Representation}

The IR is designed to make generation inspectable and locally repairable. A \texttt{SceneSpec} records required objects, object parts, materials, placement hints, cameras, screenshot plans, and spatial relations. Each required physical object can carry a collision proxy, so the harness can audit contacts and penetrations without relying on object names or natural-language descriptions alone. Materials are also explicit: \texttt{MaterialSpec} records whether image textures are required, forbidden, solid-only, or optional, and can store a resolved \texttt{texture\_source}. A \texttt{SpatialRelationSpec} separates relation semantics from checking method: horizontal support can use bounding-box contact, hinges and connectors can use attachment checks, and slanted or occluded contacts can be routed to visual verification.

Animation is represented as an extension of a validated scene rather than as a replacement for it. An \texttt{AnimationSpec} contains event ids, subject ids, target ids, frame ranges, action types, sampled frames, expected visual outcomes, and event-scoped or global \texttt{ContactConstraintSpec} entries. Contact constraints cover nonpenetration, support, touching, attachment, carry contact, and inside relations. This structure lets deterministic checks report failures with object ids, frame indices, penetration depths, overlap vectors, and tolerances, which are more actionable than free-form visual critiques.

\subsection{Progressive Scene-to-Animation Generation}

The generation process is staged because static scene errors and animation errors require different evidence. In the first stage, the original prompt is projected into a static-only IR: the scene should contain the required objects and relations, but no keyframes. The static script must pass deterministic validation and, when enabled, multi-view visual verification. Only then does the second stage add animation while preserving the validated scene geometry. This guardrail prevents an animation refiner from accidentally rebuilding the scene to hide a motion failure, and it keeps verifier feedback localized.

The end-to-end loop is deliberately simple. The planner emits an IR; the material agent resolves any requested textures; the coder writes a complete Blender script; the runtime executes it and records scene graphs, transform traces, screenshots, and previews; deterministic validators run first; optional visual and video verifiers then produce structured reports; and either deterministic repair or LLM refinement is applied until all enabled gates pass or a configured refinement budget is exhausted. In the current harness, default caps are exposed separately for deterministic, visual, and video refinement loops, and repeated identical failure signatures stop the loop early.

\subsection{Texture Resolution and Prompt-Cache-Aware Coding}

The harness resolves optional external textures before code generation when the IR says that a material benefits from image detail. A material agent uses planner-provided \texttt{texture\_policy}, \texttt{needs\_texture}, and \texttt{texture\_query} fields to search public FreeStockTextures pages, download a small candidate set, and ask the configured vision model whether any candidate is a suitable surface texture. Only approved local paths are written back as \texttt{texture\_source}; candidate counts, download counts, selected manifests, and \texttt{vision\_blocked} failures are recorded as run artifacts. If search, download, or vision approval fails, the material is marked as unavailable for image texturing and the coder falls back to base color and shader parameters. This makes texture retrieval an explicit, auditable preprocessing step rather than an implicit instruction embedded in generated Blender code.

The coding prompts were also reorganized to be compatible with provider-side prompt caching in iterative loops. The stable \texttt{ll3m\_utils} API contract is injected as a fixed system-prefix block shared by coder, refiner, user-edit, and animation-extension calls, while per-task IR and verifier evidence remain in the user payload. This engineering change does not alter the method semantics or constitute an experimental claim, but it matters operationally: repeated verifier-refiner calls reuse the same long helper contract while still receiving fresh scene and failure evidence.

\subsection{Deterministic Geometry Verification}

The deterministic validator checks both sparse semantic frames and dense temporal windows. It samples event boundaries, event midpoints, requested verifier frames, and, for short animations, all frames up to the configured audit limit. For every sampled frame, it aggregates mesh-like descendants by stable \texttt{ll3m\_id} into world-space bounding boxes. Explicit support constraints require footprint overlap and a small vertical contact gap; nonpenetration constraints measure overlap depth and axis; carry-contact constraints allow intended proximity while still detecting severe embedding. A global nonpenetration pass audits collision-enabled pairs not exempted by explicit contact relations.

These checks intentionally favor explainability over complete physical realism. Bounding boxes are conservative for concave or articulated objects, but they catch the most common failures in generated rigid animations: floating, sinking, bridge/deck crossings, conveyor penetration, detached carried objects, and final placement errors. The resulting reports are compact enough to drive automatic repair and refinement.

\begin{table}[t]
\centering
\caption{Deterministic validator predicates used by the current implementation. Tolerances are configurable in the IR; common rigid-body defaults are a $0.02$m penetration tolerance and a $0.001$m support/contact margin.}
\label{tab:validators}
\begin{tabular}{p{0.24\linewidth}p{0.34\linewidth}p{0.30\linewidth}}
\toprule
Constraint & Predicate & Reported evidence \\
\midrule
support & positive footprint overlap and bottom near support top & frame, vertical gap, footprint overlap \\
nonpenetration & bbox overlap depth below tolerance & worst frame, axis, depth, overlap vector \\
carry contact & carried object near gripper without severe embedding & frame, pairwise gap or penetration \\
inside/containment & subject center or bbox remains within container proxy & frame, escaped axes, tolerance \\
global audit & collision-enabled object pairs avoid unlicensed overlap & object pair, failing frames, worst depth \\
\bottomrule
\end{tabular}
\end{table}

\subsection{Primitive-Constrained Animation Helpers}

The primitive library exposes Blender helpers that generated scripts can call directly. Bounding-box utilities compute object tops, bottoms, centers, sizes, and bottom-to-top alignment from evaluated world-space transforms. Keyframe utilities insert location, rotation, and scale keyframes with controlled interpolation. Motion primitives then build common rigid actions from these measurements: \texttt{animate\_support\_slide}, \texttt{animate\_support\_sequence}, \texttt{animate\_pick\_place}, \texttt{animate\_attached\_carry}, \texttt{animate\_push}, \texttt{animate\_drop\_to\_support}, \texttt{animate\_rotate\_about\_axis}, and \texttt{animate\_hinge}. A dedicated gripper helper spaces two finger meshes outside the carried object's bounding box to reduce visible embedding.

The important design choice is that primitives compute contacts from executed geometry, not from the model's guessed object centers. For example, a cube on a support is positioned by aligning its aggregate bottom to the support's top plus a small margin. A pick-place motion derives source and destination rest heights from the actual support tops and the carried object's height. This makes a generated program shorter and gives the verifier less accidental geometry to repair.

\subsection{Contact-Aware Repair}

When deterministic validation finds a support or contact failure, the session can apply narrow local repairs before asking the LLM to rewrite the whole program. For horizontal bridge, deck, or platform crossings, the repair planner measures subject and support bounding boxes, narrows the support-contact window to frames where footprints overlap, and builds an outside-lift-enter-cross-exit path. For multi-support motion, sequence repair orders support windows monotonically and places keyframes at measured support centers and heights. Repair scripts rewrite only the animated subject's location curves, include descendant meshes when measuring aggregate bounds, update Blender's dependency graph before bounding-box reads, and preserve unrelated scene geometry.

\begin{table}[t]
\centering
\caption{Current deterministic repair scope. Unsupported failures are passed to the LLM refiner with deterministic and visual/video evidence.}
\label{tab:repair}
\begin{tabular}{p{0.24\linewidth}p{0.30\linewidth}p{0.34\linewidth}}
\toprule
Failure type & Diagnostic signal & Repair action \\
\midrule
horizontal support crossing & support/nonpenetration failures around bridge or deck windows & insert outside-lift-enter-cross-exit keyframes \\
multi-support transfer & ordered support constraints over road/ramp/platform-like supports & create monotonic support sequence keyframes \\
single support z error & support gap or penetration on one support ride & recalibrate z from current support top \\
conveyor ride & support object is a conveyor/belt-like ride surface & skip crossing repair and preserve ride motion \\
carry embedding & carry-contact pair shows severe penetration & report targeted evidence for primitive or LLM correction \\
\bottomrule
\end{tabular}
\end{table}

\section{System Components}

The implemented system has several components beyond the core animation method. The local Blender addon exposes execution, scene inspection, deterministic validation, screenshot rendering, animation sampling, and scene-copy commands. The Python harness records IR snapshots, generated and refined scripts, verifier reports, transform traces, selected texture assets, agent model choices, and preview summaries for each run. The LLM interface includes provider-compatibility fallbacks for JSON mode, streaming, unsupported parameters, malformed verifier JSON, and media-input capability checks. These components are not the main scientific claim, but they are necessary for making verifier-gated generation reproducible and auditable.

Material grounding is handled as an explicit preprocessing component. The planner decides whether a material should use an image texture, the material agent searches and downloads public CC0 candidates, a vision model selects only suitable surface textures, and the coder loads approved local paths into Blender materials. If any step fails, the run records the failure and falls back to non-image shader parameters. The texture pipeline is therefore isolated from the animation claims while still supporting more detailed scene appearance and preserving user constraints such as solid-only or no-image-texture prompts.

\section{Preliminary Checks and Planned Evaluation}

\subsection{Primitive Smoke Tests}

We ran local Blender smoke tests for several primitive families. These tests are implementation checks rather than a final benchmark, but they verify that the primitive code produces contact geometry consistent with the deterministic validator.

\begin{table}[t]
\centering
\caption{Preliminary primitive smoke tests. Gaps are object bottom minus support top at checked frames; values near $0.001$ correspond to the intended clearance margin.}
\label{tab:smoke}
\begin{tabular}{p{0.20\linewidth}p{0.24\linewidth}p{0.29\linewidth}p{0.15\linewidth}}
\toprule
Primitive & Scenario & Checked property & Result \\
\midrule
support slide & cube slides on table & support gap at frames 1/48/96 & 0.001 / 0.001 / 0.001 \\
pick-place & gripper transfers cube to tray & source and destination support gap & 0.001 / 0.001 \\
finger spacing & two fingers around cube & x-axis finger/cube overlap & negative overlap \\
push & pusher moves box on table & support gap at frames 1/56 & 0.001 / 0.001 \\
drop-to-support & box drops to table & final support gap at frame 96 & 0.001 \\
\bottomrule
\end{tabular}
\end{table}

The smoke tests also exposed a subtle Blender dependency-graph issue. The first bounding-box helper could read stale transforms immediately after object creation or reparenting, causing support placements to inherit unintended offsets. Updating the view layer inside the world-bounding-box helper fixed this failure and motivated using executed geometry as the primitive API boundary.

\subsection{Benchmark Protocol}

The planned benchmark will compare system variants under a fixed model configuration rather than sweeping many LLM providers. Prompt categories will include support slides, conveyor rides, table-to-tray transfers, bridge and ramp crossings, pick-carry-place interactions, pushing and pulling, dropping onto supports, hinged doors or lids, and placing objects into shallow containers. Each prompt will specify required objects, final relations, contact windows, sampled frames, and visibility requirements. Every variant will use the same prompt set, model allocation, maximum refinement budget, and static-scene verification policy unless that component is explicitly ablated.

\begin{table}[t]
\centering
\caption{Planned ablations and primary measurements.}
\label{tab:ablations}
\begin{tabular}{p{0.34\linewidth}p{0.54\linewidth}}
\toprule
Variant & Purpose \\
\midrule
direct LLM keyframes & baseline without primitive-constrained generation \\
IR without primitives & structured planning and verification, but raw keyframes \\
primitives without repair & measure first-pass benefit of geometric helper calls \\
repair without primitives & isolate deterministic correction after invalid keyframes \\
primitives plus repair & full rigid-body animation system \\
no dense-frame audit & test whether sparse sampled frames miss interpolation failures \\
no carry-contact check & test whether gripper/object errors hide behind allowed contact \\
no video verifier & measure failures not captured by deterministic geometry \\
\bottomrule
\end{tabular}
\end{table}

The main metrics will be deterministic pass rate, support violation rate, penetration rate, final relation success, carry-contact consistency, number of refinement rounds, and visual/video verifier failure labels. Because the primitive generator and deterministic validator both use bounding boxes, the final evaluation should include an independent audit: human ratings for visual plausibility and contact quality, and mesh-level or simulator-based collision checks on a subset where such checks are feasible. The benchmark should report disagreement between bounding-box validation and independent judgments.

\begin{figure}[t]
\centering
\fbox{\begin{minipage}[c][1.25in][c]{0.92\linewidth}
\centering
\textbf{Placeholder: failure taxonomy and metrics figure.}\\
Rows: support, penetration, carry contact, final relation, visibility.\\
Columns: direct keyframes, primitives, repair, full system.
\end{minipage}}
\caption{Evaluation placeholder. The final figure should visualize which failure modes are caught by deterministic checks, video verification, and repair.}
\label{fig:taxonomy}
\end{figure}

\section{Limitations}

The current system is kinematic and verification-oriented, not a replacement for physical simulation. The primitives create validator-aligned rigid motions by construction, but they do not model forces, friction, mass, torque, or collision response. Bounding-box checks can be overly strict for concave objects and insufficient for fine-grained mesh intersections. Slanted ramps, articulated multi-link arms, deformable objects, fluids, and cloth require richer geometry or simulation support than the rigid primitive set currently provides. Finally, video-language verification is useful for semantic and temporal judgments, but it can miss short contact failures; it should supplement deterministic checks rather than replace them.

\section{Conclusion}

We presented a verification-oriented harness for LLM-generated rigid-body animation in Blender. The system represents scene and motion intent in a structured IR, generates animation only after validating a static scene baseline, checks contact and penetration constraints with deterministic geometry, and repairs common failures from measured scene-graph evidence. A new library of primitive-constrained animation helpers moves repeated support, carry, push, drop, and hinge motions out of unconstrained LLM keyframe text and into reusable geometry-aware operations. The next step is a controlled benchmark that quantifies how much primitive-constrained generation and deterministic repair improve validator-measured contact validity and refinement efficiency, with independent audits to test whether those gains correspond to visual and physical plausibility.

\bibliography{references}
\bibliographystyle{iclr2026_conference}

\end{document}
